\documentclass[a4paper]{article}

% Basic packages
\usepackage{fontspec}\setmainfont{Noto Serif CJK SC}
\usepackage{amsmath, amssymb}
\usepackage{graphicx}
\usepackage[margin=2.5cm]{geometry}
\usepackage[most]{tcolorbox}
\usepackage{etoolbox}
\usepackage{listings}
\usepackage{hyperref}
\usepackage{booktabs}
\usepackage{subcaption}
\usepackage{float}
\usepackage{tikz}
\IfFileExists{pgfplots.sty}{
    \usepackage{pgfplots}
    \pgfplotsset{compat=1.18}
}{}

% Highlight boxes
\newtcolorbox{knowledgebox}[1]{
    enhanced, colback=blue!5!white, colframe=blue!75!black, colbacktitle=blue!75!black,
    coltitle=white, fonttitle=\bfseries, title=#1, attach boxed title to top left={yshift=-2mm, xshift=2mm},
    boxrule=1pt, sharp corners
}

\newtcolorbox{importantbox}[1]{
    enhanced, colback=yellow!10!white, colframe=yellow!80!black, colbacktitle=yellow!80!black,
    coltitle=black, fonttitle=\bfseries, title=#1, sharp corners
}

\newtcolorbox{warningbox}[1]{
    enhanced, colback=red!5!white, colframe=red!75!black, colbacktitle=red!75!black,
    coltitle=white, fonttitle=\bfseries, title=#1, sharp corners
}

% Listings
\lstset{
    language=Python,
    basicstyle=\ttfamily\small,
    keywordstyle=\color{blue},
    stringstyle=\color{red!60!black},
    commentstyle=\color{green!60!black},
    breaklines=true,
    frame=single,
    numbers=left,
    numberstyle=\tiny\color{gray},
    captionpos=b,
    extendedchars=false
}

% Front-page metadata
\newcommand{\notetitle}{基本共射放大电路的动态分析}
\newcommand{\noteauthors}{五道口纳什 \& Codex}
\newcommand{\notedate}{\today}
\newcommand{\videochannel}{上海交通大学 郑益慧}
\newcommand{\videopublishdate}{[在此填写发布日期]}
\newcommand{\videoduration}{约 46 分钟}
\newcommand{\videourl}{}
\newcommand{\videocoverpath}{cover.jpg}

\begin{document}

\begin{titlepage}
\centering
{\Large 课程笔记\par}
\vspace{1.2cm}
{\huge\bfseries \notetitle\par}
\vspace{0.8cm}
{\large \noteauthors\par}
\vspace{0.3cm}
{\large \notedate\par}
\vspace{1.2cm}

\ifdefempty{\videocoverpath}{
{\small 请在生成笔记时填入视频封面图的本地路径。\par}
}{
\includegraphics[width=0.82\textwidth,height=0.45\textheight,keepaspectratio]{\videocoverpath}\par
}

\vfill
\begin{tcolorbox}[width=0.9\textwidth, colback=black!2!white, colframe=black!60, sharp corners]
\textbf{视频作者/频道}：\videochannel\par
\textbf{发布日期}：\videopublishdate\par
\textbf{视频时长}：\videoduration\par
\textbf{视频链接}：
\ifdefempty{\videourl}{
未填写
}{
\href{\videourl}{\nolinkurl{\videourl}}
}
\end{tcolorbox}
\end{titlepage}

\tableofcontents
\newpage

%% --- 正文内容开始 --- %%
\section{问题背景：从静态工作点进入动态分析}
本讲讨论基本共射放大电路。已知电源、电阻、晶体管静态参数和偏置条件后，目标是求输入输出电压关系、放大倍数以及输入、输出电阻。核心方法是先求静态工作点，再建立交流等效模型。
\begin{importantbox}{动态分析的顺序}
先画直流通路求 $I_{BQ},I_{CQ},U_{CEQ}$，再画交流通路并用小信号模型求 $r_{be}$、$A_u$、$R_i$、$R_o$。
\end{importantbox}
\subsection{本章小结}
动态问题不能跳过静态工作点，因为小信号参数依赖工作点。

\section{静态工作点与电位分析法}
直流分析时，耦合电容和旁路电容视为断路。沿电路逐元件追踪电位，可以稳定地列出 KVL。例如基极回路可写为
\[
V_{CC}=U_{BQ}+I_{BQ}R_B,\qquad I_{CQ}=\beta I_{BQ}.
\]
其中 $V_{CC}$ 是电源电压，$U_{BQ}$ 是基极静态电位，$R_B$ 是偏置电阻，$I_{BQ}$ 和 $I_{CQ}$ 分别是基极、集电极静态电流。
\begin{warningbox}{常见错误}
不要把并非同一支路的电阻与 $I_{BQ}$ 强行相乘。应逐段判断电流，再用每个元件两端的电位差列式。
\end{warningbox}
\subsection{本章小结}
电位分析法把复杂回路拆成一段一段的电位上升和下降，适合检查偏置计算。

\section{交流通路与小信号等效}
交流分析时，直流电源置零，耦合与旁路电容在中频段视为短路。晶体管替换为小信号等效模型，输入端看到 $r_{be}$，受控电流近似为 $\beta i_b$。这一步的意义是把非线性器件在工作点附近线性化。
\[
A_u=\frac{u_o}{u_{is}},\qquad R_i=\frac{u_i}{i_i},\qquad R_o=\frac{u_o}{i_o}.
\]
\subsection{本章小结}
交流通路只保留信号变化；静态偏置仍然决定小信号参数，但不再以独立直流源形式出现在交流等效图中。

\section{负载、放大倍数与输出电阻}
接入负载 $R_L$ 后，集电极等效负载通常为 $R_C\parallel R_L$，因此电压增益的绝对值随负载变化。输出电阻用端口测试法求：从输出端看入放大电路，独立输入信号置零，再计算端口电压与电流之比。
\begin{importantbox}{设计取舍}
若目标是输出电压，较小的输出电阻有利于减小负载造成的衰减；若目标、级间匹配或功率传输不同，最佳取值也会不同。
\end{importantbox}
\subsection{本章小结}
$R_L$ 不只是附加元件，它会改变交流等效负载，从而改变增益和工作范围。

\section{总结与延伸}
本讲的主线是“静态决定线性化位置，交流模型决定增量关系”。面对基本共射放大电路，可按以下清单操作：先断开电容求静态工作点；再将电源置零、短接电容建立交流通路；替换晶体管小信号模型；最后分别求增益、输入电阻和输出电阻。视频后段还强调，输出电阻大小应服从应用目标，不能脱离负载和信号类型孤立判断。


% Replace this block with the generated notes.
% Fill the metadata block above, especially \videocoverpath, so the front page
% contains the original video cover image.
% Keep images outside knowledgebox/importantbox/warningbox.
% For any figure that comes from a video frame or a crop of a frame,
% include the source time interval as a bottom footnote on the same page.
% A stable pattern is:
% \begin{figure}[H]
% \centering
% \includegraphics[width=\textwidth]{...}
% \caption{... \protect\footnotemark}
% \end{figure}
% \footnotetext{视频画面时间区间：00:12:31--00:12:46。若为裁剪图，时间区间仍对应原视频画面。}
% End the document with a final section such as:
% \section{总结与延伸}
% covering the speaker's closing discussion and your own synthesis.

%% --- 正文内容结束 --- %%

\end{document}
